\section{Support System Integration and Verification}
\label{section:integrate-and-verify-system}

Verification is planned. OOSEM specifies stimulus, conditions, expected response and verdict and treats the environment as an enabling system \citep[Secs.~17.3.8--17.3.9]{friedenthal2014practical}. The Handbook places verification approach and AIV planning before execution/status reports \citep[Secs.~I.3.1--I.3.6]{aguado2021cubesatse}. No framework procedure has run.

The unit under verification is an identified implementation/core. Declare service/tool versions, legal access, CPU solvers, controlled fixtures, backups and clean independent replay. Qualified reviewers adjudicate references. Fix numerical rules before execution. Missing or unjustified criteria leave verdicts undetermined. Table~\ref{tab:framework-verification-increments} covers all forty V IDs. Each verifies its same-numbered FR. Integrate prerequisites in order and repeat affected checks after changes.
{\small
\setlength{\tabcolsep}{4pt}
\setlength{\parskip}{0pt}
\setlength{\LTpre}{0.4\baselineskip}
\setlength{\LTpost}{0.4\baselineskip}
\renewcommand{\arraystretch}{1.06}
\begin{longtable}{@{}>{\raggedright\arraybackslash}p{\dimexpr0.29\linewidth-0.58\tabcolsep\relax}>{\raggedright\arraybackslash}p{\dimexpr0.71\linewidth-1.42\tabcolsep\relax}@{}}
\caption{Planned increments, prerequisite fixtures and expected integration evidence.}\label{tab:framework-verification-increments}\\
\toprule
\textbf{Increment / V coverage} & \textbf{Prerequisite and expected response} \\
\midrule
\endfirsthead
\caption[]{Planned increments, prerequisite fixtures and expected integration evidence. (continued).}\\
\toprule
\textbf{Increment / V coverage} & \textbf{Prerequisite and expected response} \\
\midrule
\endhead
\bottomrule
\endfoot
IN-01 / V01--03, V15 & Core schema/guidance and qualitative/quantitative fixtures. Round-trip fields, source methods and guided case instantiation. \\
IN-02 / V04--09, V11 & Reference graph, corrupt transfers and interrupted changes. Check impact, exact history, rejection, failed gates and stale/critical status. \\
IN-03 / V12--14, V17, V26--27 & Sourced components/quotes and expert RAM reference. Check BOM/arithmetic, derivation, comparison and failure transformation with unknowns. \\
IN-04 / V16, V18--20, V23--24 & Pinned dynamic/tool/battery/controller SIL and policy fixtures. Check shortage/protection, mappings, coupling, raw results, timing and limitations. \\
IN-05 / V21--22, V25, V31--32 & Distinct calibration/validation, domain and fault/degradation/noise/clock fixtures. Check residuals, updates, alternatives, raw channels and masks. \\
IN-06 / V33--36 & Pinned cFS/F-Prime dictionaries, sequence and target references. Check mappings, guards/effects/recovery and resource/compatibility gaps. \\
IN-07 / V28--30 & Method-specific and incomplete/unsafe proposals with technician/proxy review. Check DUT/equipment, capture and blocked readiness without hazardous execution. \\
IN-08 / V10, V37--40 & Released baseline, clean replay, reviewer tasks and both frozen-core cases. Check reconstruction/views and record actual effort/cost/core changes. \\
\end{longtable}
}

V18 requires a successful reference exchange and semantic checks for every electrical, thermal, EMC and CAD profile required by the frozen minimum. A documented development gap leaves its obligation not verified or undetermined. A narrower minimum needs a reviewed scope disposition. V28 includes quantitative and qualitative fixtures with a rationale for inapplicable fields.

The registry defines all forty procedures once: subject/revision, fixture, conditions, steps, expected response, acceptance observation, evidence, reviewer and unresolved decisions. Retain raw inputs/outputs/logs, environment/model/criterion versions, masks, clock records and findings. Current status is not executed. After execution, pass, fail or undetermined refers to the named subject and fixed rule.

Framework verification checks FR-01--FR-40. Usefulness validation observes MOEs through stakeholder tasks. Case EPS acceptance checks that EPS's requirements. Native model syntax/semantic checks establish model consistency within their tool scope. These are separate. Later physical campaigns need calibrated equipment, qualified clocks/delays, safe fault limits, isolation, abort/recovery and facility authorization. Virtual fixtures can verify blocking without hazardous work.

The project proposal preceded the expanded need analysis. Its objectives, deliverables, indicators and deadlines form a separate user-supplied research baseline. Results will assess them separately from framework verdicts and case EPS evidence. They are not populated here. The chapter ends at specified design. Stakeholder review, implementation and empirical verification remain pending.
